\markdownRendererDocumentBegin
\markdownRendererSectionBegin
\markdownRendererHeadingOne{Limitations, discussion and conclusion}\markdownRendererInterblockSeparator
{}The exact theorem is restricted to a declared linear-access class and is not an unrestricted representation or time lower bound. Parity families are deliberately controlled. P11D shows that stronger decoder bias materially reduces compilation advantage. P11G tests one finite nonlinear ensemble. The digits study is a ten-responsibility, one-dataset boundary rather than an external generalization test: it supports only 3/10 responsibilities under linear access and 5/10 under each stronger access class, and it fails the 8/10 family gate without retuning. Controlled operation counts do not replace end-to-end latency, memory traffic, training cost or energy. Future-query laws assume the declared workload model. Broad real-system superiority still requires matched same-information experiments in real LLM/procedural and formal/search or long-horizon settings.\markdownRendererInterblockSeparator
{}The experiments change the interpretation of task-conditioned construction. The strongest story is not that a compiler creates information or universally beats a universal representation. The decisive variable is \markdownRendererStrongEmphasis{who performs structural search}. A query-conditioned compiler performs some search before the downstream learner sees state. A sparse decoder can recover part of that work; a deterministic finite tree ensemble is another search mechanism but remains inefficient under the registered high-dimensional low-sample envelope.\markdownRendererInterblockSeparator
{}Aggressive specialization can make one task cheap while destroying immediate support for future tasks. Retaining raw state, caching compilations or materializing wider state are different allocations of computation and memory across time.\markdownRendererInterblockSeparator
{}P11 establishes a controlled theory/systems result: fixed linear-accessible state scales with query-family rank; task-conditioned construction can expose smaller task-facing state; a hostile sparse decoder buys part of the controlled gain back; and a width-conditioned \markdownRendererCodeSpan{r=7} result survives the registered pooled attack. The adverse digits result prevents promotion to family-scale compilation support. \markdownRendererCodeSpan{P11\markdownRendererUnderscore{}ACTIVE\markdownRendererUnderscore{}CLAIM\markdownRendererUnderscore{}AUTHORITY\markdownRendererUnderscore{}V2.json} binds these disjoint positive and negative scopes without converting the latter into support.\markdownRendererInterblockSeparator
{}\markdownRendererBlockQuoteBegin
\markdownRendererStrongEmphasis{State construction, decoder search, and future recoverability are coupled resource choices. Moving structure into state can reduce downstream discovery work, but the benefit shrinks as the downstream access mechanism becomes better at discovering that structure, and specialization creates option debt unless recoverability is retained.}
\markdownRendererBlockQuoteEnd \markdownRendererInterblockSeparator
{}The controlled claim does not authorize a universal nonlinear lower bound, transformer scaling law, free preprocessing or real-agent superiority.
\markdownRendererSectionEnd \markdownRendererDocumentEnd